\section{EA Test 5: valley flooding}

Test~5 of the Environment Agency benchmark \cite{NeelzPender2013} follows
the flood from a failed embankment dam 17~km down a real valley, which
falls at about 1:100 in its upper part and 1:1000 in its lower part, to a
pond at its foot (Figure~\ref{fig:ea5_setup}). The inflow peaks at
3000~m$^3$/s between 10 and 20~min and stops at 100~min, entering along a
260~m line across the head of the valley; the modelled area (19.1~km$^2$)
is the EA's outline, closed on all sides, on a dry bed. Manning's
$n = 0.04$, the specified resolution is 50~m (here triangles of at most
2500~m$^2$, about 11,900 in all), and the run lasts 30~h. Ground levels
are sampled at the triangle centroids from the 10~m DEM.

\begin{figure}[h]
\begin{center}
\includegraphics[width=0.95\textwidth]{setup.png}
\end{center}
\caption{EA Test 5: the valley, the modelled area, the inflow line and the
seven output points, and the inflow.}
\label{fig:ea5_setup}
\end{figure}

In the report (Section~4.6.4) the full shallow-water packages took almost
3~h to reach point~5, agreeing within about 30~min, and their final levels
in the pond at point~5 lay within 0.4~m of each other. Peak speeds lay in
1.5--2.2~m/s at point~4 and 1.2--1.6~m/s at point~7. \anuga{}~1.1 was the
exception in 2009: it reached point~5 about 1.5~h late, its final level
there was about 0.5~m below the others, and it showed spurious high speeds
in very shallow flow. The current \anuga{} reaches point~5 at 3.3~h and
settles there at 148.56~m, where TUFLOW's five current solvers give
2.8--3.3~h and 148.6--148.9~m \cite{BMT2017}, and its peak speeds at
points 4 and 7 lie within the report's ranges. The answer is converged
in the mesh: with triangles of 625 and 100~m$^2$ (about 47,500 and
293,000) the arrival at point~5 moves by under 5~min, the final levels by
at most 0.11~m (at point~7, rising towards TUFLOW's) and the peak speeds
by under 0.1~m/s. The checks
(Figure~\ref{fig:ea5_checks}) are the volume balance, the order in which
the flood reaches the points, these figures from the report, and a stored
baseline per flow algorithm.

\begin{figure}
\begin{center}
\includegraphics[width=0.95\textwidth]{stage_plot.png}
\end{center}
\caption{EA Test 5: water level at the seven output points.}
\end{figure}

\begin{figure}
\begin{center}
\includegraphics[width=0.95\textwidth]{speed_plot.png}
\end{center}
\caption{EA Test 5: speed at the seven output points.}
\end{figure}

\begin{figure}
\begin{center}
\includegraphics[width=0.9\textwidth]{checks.png}
\end{center}
\caption{EA Test 5: the checks and their outcome.}
\label{fig:ea5_checks}
\end{figure}
